% Personal preparation companion. Main report reordered on 02/10/2026; numbers unchanged.
\documentclass[10pt,a4paper]{article}
\usepackage{fontspec}
\setmainfont{texgyretermes}[Extension=.otf,UprightFont=*-regular,BoldFont=*-bold,ItalicFont=*-italic,BoldItalicFont=*-bolditalic]
\usepackage[margin=1.9cm,headheight=14pt]{geometry}
\usepackage{amsmath,amssymb,booktabs,longtable,array,ragged2e,enumitem,microtype,xcolor,graphicx,fancyhdr,tikz}
\usetikzlibrary{arrows.meta,positioning}
\usepackage[hidelinks,unicode]{hyperref}
\definecolor{ink}{HTML}{183A45}\definecolor{teal}{HTML}{087F7A}\definecolor{pale}{HTML}{EAF4F2}\definecolor{gray}{HTML}{56656B}
\newcolumntype{P}[1]{>{\RaggedRight\arraybackslash}p{#1}}
\setlength{\parindent}{0pt}\setlength{\parskip}{6pt}\setlength{\emergencystretch}{2em}
\setlength{\tabcolsep}{4pt}\renewcommand{\arraystretch}{1.15}
\setlist{nosep,leftmargin=1.4em,topsep=4pt,itemsep=3pt}
\pagestyle{fancy}\fancyhf{}\fancyhead[L]{\small\color{gray}GHI CHÚ CHUẨN BỊ TRÌNH BÀY}\fancyhead[R]{\small\color{gray}03/10/2026}\fancyfoot[C]{\small\thepage}\renewcommand{\headrulewidth}{0.3pt}
\newcommand{\key}[1]{\par\smallskip\noindent\colorbox{pale}{\parbox{\dimexpr\linewidth-2\fboxsep}{#1}}\par\smallskip}
\newcommand{\say}[1]{\par\textbf{Có thể trình bày:}\ \textit{#1}\par}
\newcommand{\ask}[2]{\par\textbf{Nếu được hỏi: #1}\par #2\par}
\newcommand{\where}[1]{\par{\small\color{gray}Đọc cùng report chính: #1.}\par}
\newcommand{\transition}[1]{\par\textbf{Chuyển ý:}\ #1\par}
\newcommand{\example}{\textbf{Ví dụ minh họa, không phải kết quả benchmark.} }
\newcommand{\guidepage}[1]{\clearpage\section{#1}}
\newcommand{\layer}[2]{\clearpage\noindent{\color{teal}\small\bfseries #1}\par\vspace{2pt}{\Large\bfseries\color{ink}#2}\par\vspace{6pt}\hrule\vspace{8pt}}
\hypersetup{pdftitle={Chuẩn bị trình bày nghiên cứu: related works, kiến trúc và benchmark},pdfauthor={}}
\input{preparation_evidence.tex}
\begin{document}
\begin{center}
{\LARGE\bfseries\color{ink}Chuẩn bị thuyết trình}\\[5pt]
{\large Related works → kiến trúc → benchmark}\\[3pt]
{\small Buổi 03/10/2026 · Đọc cùng report rút gọn · Không có thực nghiệm mới}
\end{center}
\key{Mạch nói 12–15 phút: 4 phút về papers/metrics; 4 phút về kiến trúc; 4–6 phút về kết quả. Phụ lục để mở khi cần giải thích A/B/C.}

\section{Related works và lý do chọn metrics}
\say{Em đã cập nhật related works sang các công trình trong ba năm gần đây, khóa khảo sát ở 21/09/2023–21/09/2026. DLS và RDG giữ làm mốc nền; AnotherMe có số tạp chí 2024 nhưng online năm 2023 ngoài cửa sổ nên ghi là ngoại lệ. Từ bảng metrics, em thấy các paper đang chấm những mục tiêu khác nhau, vì vậy cần bộ đo chung cho dịch vụ của mình.}

\begin{longtable}{@{}P{4.1cm}P{5.4cm}P{6.1cm}@{}}
\toprule Paper tiêu biểu & Metrics gốc & Điều cần rút ra\\\midrule
TransProtect 2024 & EIE; chênh chi phí hành trình & Sai số suy luận gần nhiệm vụ privacy; chi phí hành trình chưa đo đúng top-5 POI.\\
Semantic / Fake queries 2026 & ASR, DER hoặc số đường khó phân biệt; thời gian sinh & ASR là thành công ẩn danh theo định nghĩa paper; không tự đổi thành tỷ lệ attacker đoán đúng.\\
DP-FETC 2025 & MI, JSD; bảo đảm DP & Bảo toàn phân bố quỹ đạo chưa bảo đảm dịch vụ POI tốt hoặc endpoint khó suy.\\
CPCROK / PRISM 2025 & Khớp quỹ đạo, số giả; LPI, TDD, DC & Khác tác vụ hoặc thiếu định nghĩa được xác minh; không dùng tên metric để so ngang.\\
SoK 2024 & Kiểm tra tấn công và utility theo tác vụ & Hỗ trợ cách đánh giá trên cùng nhiệm vụ, không phải một model đối chứng.\\
\bottomrule
\end{longtable}
Bảng đầy đủ 12 nguồn gần đây nằm ở mục 1 của report; phần chưa đủ nguồn được ghi CX.

\textbf{Ba câu hỏi dẫn tới bộ đo đề xuất:}
\begin{itemize}
\item \textbf{Attacker suy đúng tới đâu?} Hit50/100/200, MAE, median, p90. Hit100 là tỷ lệ suy ra vị trí cách đáp án ≤100 m; thấp hơn tốt hơn cho privacy. Đọc Hit cùng sai số để tránh MAE lớn che một số lần lộ chính xác.
\item \textbf{Người dùng nhận được gì?} Recall@5, trả đủ và tăng khoảng cách đường $\Delta D$. So với top-5 chuẩn dùng GPS thật, cùng trạng thái khả dụng. Recall từng ca cần ≥90\%.
\item \textbf{Trả giá bao nhiêu?} Byte yêu cầu + phản hồi mỗi sự kiện thật, cả query chèn và tải ngoài chuyến; latency p50/p95 khi có. Hiện chưa đo latency đầu-cuối, chưa tính Q thực nghiệm.
\end{itemize}
\say{Em không đề xuất phát minh MAE hay Recall. Em đề xuất cách chọn và vận hành hóa các chỉ số để một tập dummy và một vị trí thay thế đều được chấm trên cùng đáp án, cùng dịch vụ và chi phí đầy đủ.}
\transition{Sau khi xác định thế nào là bảo vệ tốt mà vẫn phục vụ được, em thiết kế luồng truy vấn như sau.}

\clearpage\section{Kiến trúc mô hình: Geo-I và các lớp cơ chế}
\begin{center}\includegraphics[width=\linewidth]{figures/architecture_model.pdf}\end{center}
\say{Nền tảng của mô hình là Geo-I. Từ vị trí thật, em tạo neo riêng tư; các bước sau dùng lịch sử neo cùng mạng đường và POI công khai để sinh một tập điểm giả. Em bổ sung xử lý biên cho điểm đầu/cuối, thay vì coi nhiễu từng vị trí là đủ.}
\begin{enumerate}
\item \textbf{Geo-I / REM — S1:} lấy neo bằng cơ chế mũ trên tập vị trí đường công khai, dùng khoảng cách Euclid. Đây là nền tảng bảo vệ từng vị trí; không chỉ thêm nhiễu rồi snap tùy ý.
\item \textbf{Tái dùng neo, ngân sách và pacing — S2/S3:} test có nhiễu để giữ hoặc tạo neo; tính cả hai loại chi phí. Giảm báo nhiều mẫu nhiễu độc lập khi dừng và kiểm soát composition theo phiên. Hết ngân sách thì tiếp tục hậu xử lý, không đọc GPS mới. Pacing là giãn đọc GPS, chưa che giờ bắt đầu/kết thúc chuyến.
\item \textbf{Belief + road-aware + POI-aware — S2/S3 và utility:} ước lượng vị trí/chuyển động từ neo; ràng buộc dummy đi được trên làn có hướng, đúng luật rẽ và thời gian; chọn K điểm phủ các POI bổ sung nhau. Switching thêm mode dừng/đi; exchange/slack là biến thể bộ chọn. Selector không đọc tốc độ thật hay đích tương lai.
\item \textbf{BR-Boundary — S9/S10:} bỏ h giây đầu \emph{trước lõi} để điểm đầu không ảnh hưởng neo; giữ bản tin trong buffer Δ giây \emph{sau lõi}, kết thúc thì hủy phần chưa phát. Không cần đoán trước 60 giây cuối.
\end{enumerate}
\key{Lõi S1–S3: Geo-I + xử lý thời gian + mạng đường/POI. Mở rộng S9/S10: bỏ đầu trước lõi và buffer đuôi sau lõi. Các lớp boundary đã có kiểm tra code/tích hợp, chưa có benchmark toàn bộ mô hình với dịch vụ.}
\ask{Road-aware có tự tăng bảo đảm Geo-I không?}{Không. Nó làm dummy khả thi trên đường và hỗ trợ dịch vụ. Cận riêng tư đến từ cơ chế neo và cách cộng ngân sách; hậu xử lý chỉ dùng dữ liệu đã bảo vệ giữ cận đó. Kháng attacker biết đường vẫn phải đo. K điểm giả không tự là k-anonymity.}
\transition{Cần phân biệt mô hình vừa mô tả với nhánh truy vấn công khai tạo ra các bảng kết quả tiếp theo.}

\clearpage\section{Benchmark: ghi đúng mô hình tạo ra kết quả}
\key{Bảng 30/67 dưới đây thuộc nhánh kế hoạch truy vấn công khai, không dùng neo Geo-I. Không đọc Hit100=0 thành kết quả của lõi Geo-I hoặc của Geo-I + BR-Boundary.}
\textbf{Giới thiệu đối chứng một lần:} DLS, RDG, TransProtect*, Semantic*, Fake-query* là năm adapter online. TransProtect dùng Markov thay Transformer; Semantic dùng predictor thực nghiệm thay LSTM. AnotherMe gốc là online, adapter VTGA hiện có đọc cả chuyến nên chỉ làm tham chiếu offline.
\PrivacyEvidence
\say{Hit100 của nhánh công khai bằng 0 trong bộ attacker đã thử, trong khi các đối chứng vẫn bị suy đúng gần. Nhưng S1–S3 là bốn nhóm và client theo sự kiện; S9–S10 là 32 nhóm và client theo lịch. Em không gộp hai phạm vi này thành một phép thử duy nhất.}

\textbf{Bảng endpoint trên 32 nhóm.} Privacy là Hit100 / MAE mét; utility tại 80\% khả dụng; byte mỗi sự kiện thật.
\EndpointEvidence
\begin{itemize}
\item \textbf{Vì sao có sức thuyết phục?} Raw có hit ở cả S10.A/B; CI 95\% của chênh lệch Hit100 so với từng adapter nằm dưới 0 ở S9 và S10. Bootstrap theo nhóm; chưa hiệu chỉnh nhiều so sánh.
\item \textbf{Đánh đổi 30/67:} bản 30 đạt Recall 95,00\%, 14/14 ca ≥90\% ở p=0,8; stress p=0,95 còn 93,01\%, 13/14 ca. Bản 67 đạt 100\%, 14/14 ca ở các mức đã thử, nhưng tốn khoảng 6.879 thay vì 3.071 byte/sự kiện.
\item \textbf{Giá của lịch:} 4,48 lần byte so cùng kế hoạch chỉ refresh khi hoạt động. Đổi lại giờ bắt đầu/kết thúc không điều khiển lịch tải. Byte chưa gồm HTTP/TLS.
\item \textbf{AnotherMe:} số privacy chỉ trên tập con sinh thành công; utility giữ lỗi bằng 0. Không xếp như cùng mẫu số với năm adapter online.
\end{itemize}
\say{Kết quả hỗ trợ nhánh công khai trên năm scenario hiện tại; chưa xác nhận toàn bộ kiến trúc Geo-I và boundary vừa mô tả. Privacy được cải thiện với chi phí truyền tăng; em chưa kết luận vượt các paper nguyên bản hoặc tốt hơn ở cùng ngân sách.}
\ask{Hit100=0 có nghĩa bảo vệ 100\% không?}{Không. Đây là không trúng trong các mẫu và attacker đã thử. Còn prior, kênh mạng ngoài phạm vi, thành phố khác và S1–S3 trên client theo lịch cần được kiểm chứng. Bản 67 vẫn có tám POI ngoài độ phủ tĩnh.}

\clearpage
\input{scenario_appendix.tex}
\end{document}
